\documentclass{article}
\usepackage[utf8]{inputenc}
\usepackage[T1]{fontenc}
\usepackage[french]{babel}
\usepackage{graphicx}
\usepackage{hyperref}
\usepackage{tikz}

\title{Rapport Technique : Application Traveling}
\author{BERGON Paul, VERAN Louis}
\date{\today}

\begin{document}

\maketitle
\tableofcontents
\newpage

\section{Introduction et description du projet}
L'application \textbf{Traveling} est une application mobile Android de voyage composée de deux entités majeures connectées par une passerelle :
\begin{itemize}
    \item \textbf{TravelShare} : Un espace de partage et de découverte de photos de voyages (mode public/anonyme ou connecté).
    \item \textbf{TravelPath} : Un générateur intelligent de parcours de visite basé sur les préférences de l'utilisateur.
\end{itemize}

\section{Architecture globale de l'application}

\begin{tikzpicture}[node distance=2cm, align=center]
  \node[draw, rounded corners, minimum width=3cm, minimum height=1cm] (app) {Application Android (UI)};
  \node[draw, below left of=app, node distance=4cm, fill=blue!10] (firebase) {Firebase (Auth, Firestore, FCM)};
  \node[draw, below right of=app, node distance=4cm, fill=cyan!10] (cloudinary) {Cloudinary (Médias)};

  \draw[<->] (app) -- node[left] {Données \& Sync} (firebase);
  \draw[<->] (app) -- node[right] {Upload \& CDN URL} (cloudinary);
\end{tikzpicture}

\section{Description détaillée des fonctionnalités}

\subsection{Phase 3 : Sous-système TravelShare}
Le module \textbf{TravelShare} a été développé avec une interface utilisateur immersive de type "flux photographique" (similaire à Instagram). Au cœur de cette vue se trouve le \texttt{HomeFragment}, qui implémente un \texttt{RecyclerView} optimisé.
Les cartes présentent l'avatar du posteur, la localisation, la photo (servie avec une mise en cache via la librairie Glide), le titre, la description ainsi que le nombre de "\textit{likes}". Le routing entre les écrans passe par le composant racine \texttt{NavGraph}.

\subsubsection{Composant : \texttt{PhotoDetailActivity} (Détail Dynamique)}
Une page de détail complète, \texttt{PhotoDetailActivity}, a été conçue pour offrir une immersion totale dans chaque cliché partagé.
\begin{itemize}
    \item \textbf{Dynamisme et Parcelable} : Contrairement à une vue statique, cet écran est 100\% piloté par les données. La classe \texttt{Photo} implémente l'interface \texttt{Parcelable} pour permettre le transfert d'objets complexes (incluant les coordonnées GPS, les métadonnées d'auteur et les listes de tags) via les \texttt{Intents} Android.
    \item \textbf{Structure du Layout} : Le fichier \texttt{activity\_photo\_detail.xml} repose sur un \texttt{NestedScrollView}. Il intègre une bannière de 250dp, une section auteur riche, un encart call-to-action (CTA) raffiné (fond blanc, bordure \textit{teal} \#009688, icône en bulle), et des sections sociales étendues.
    \item \textbf{Système de Commentaires} : Une section interactive permet de visualiser les avis de la communauté via le \texttt{CommentAdapter}. L'utilisateur peut saisir un nouveau commentaire, déclenchant un ajout asynchrone en haut de liste et une mise à jour dynamique du compteur d'engagement (\textit{feedback} immédiat).
    \item \textbf{Recommandations Intelligentes} : Sous les commentaires, un carrousel horizontal affiche des "Photos similaires" (ex: Arc de Triomphe, Louvre) géré par \texttt{SimilarPhotoAdapter}. Ce module favorise la rétention utilisateur en permettant une navigation fluide entre lieux connexes de la capitale.
    \item \textbf{Harmonisation Chromatique} : Toutes les couleurs ont été alignées sur les \textit{Design Tokens} officiels (\texttt{@color/primary}, \texttt{@color/secondary}, \texttt{@color/muted}). L'encart CTA bénéficie d'un arrière-plan en dégradé subtil (5\% d'opacité) avec une bordure à 20\%, offrant un rendu moderne et premium conforme aux spécifications Tailwind.
    \item \textbf{Interactions et Connectivité} :
    \begin{itemize}
        \item \textbf{Navigation \& GPS} : Intégration de liens profonds vers Google Maps via des \texttt{Intent.ACTION\_VIEW} utilisant les coordonnées \texttt{lat/lng} stockées dans l'objet.
        \item \textbf{Engagement} : Système de \textit{toggles} interactifs pour le "Like" et le "Bookmark" avec retour visuel immédiat (changement de couleur et d'icône Lucide).
        \item \textbf{Passerelle TravelPath} : Un bouton "Créer un parcours" permet de basculer vers le module de planification en injectant automatiquement la localisation de la photo comme point d'intérêt.
        \item \textbf{Gestion de l'Encoche} : Utilisation de \texttt{ViewCompat.setOnApplyWindowInsetsListener} pour appliquer dynamiquement un \textit{padding} supérieur, évitant que la barre de titre ne soit masquée par la caméra frontale.
    \end{itemize}
    \item \textbf{Difficultés \& Solutions} : Le rendu dynamique des "Tags" sous forme de \textit{Chips} Material a nécessité l'utilisation d'un \texttt{ChipGroup} peuplé programmatiquement en Java, afin de garantir une adaptation parfaite à la longueur variable des thématiques (ex: "Monument", "Coucher de soleil"). 
\end{itemize}

\subsection{Publication de Contenu}
L'activité \texttt{PublishPhotoActivity} a été conçue pour offrir un flux de création intuitif et riche :
\begin{itemize}
    \item \textbf{Sélection de Photo} : Zone d'upload dynamique avec bordure en pointillés (\textit{dotted border}). Gestion de la galerie via \texttt{ActivityResultLauncher}.
    \item \textbf{Métadonnées Enrichies} : Saisie du titre, de la description, de la date (via \texttt{DatePickerDialog}) et du moment de la journée (Menu Popup).
    \item \textbf{Note Audio (Stub)} : Emplacement prévu pour l'ajout de commentaires vocaux, renforçant l'aspect narratif du voyage.
    \item \textbf{Indexation par Tags} : Système de \texttt{ChipGroup} permettant l'ajout manuel et la suggestion assistée par "IA" (basée sur des mots-clés du titre).
    \item \textbf{Harmonisation Chromatique} : Comme pour le reste de l'application, l'activité de publication utilise désormais exclusivement les \textit{Design Tokens} (\texttt{@color/primary}, \texttt{@color/secondary}, \texttt{@color/muted\_foreground}) pour les boutons, icônes et bordures, garantissant une cohérence visuelle totale.
    \item \textbf{Simulation de Persistance} : Les photos publiées sont injectées dans le \texttt{MockDataProvider} pour une mise à jour immédiate du profil utilisateur.
\end{itemize}

\subsection{Stabilisation Technique}
Une phase de refactorisation profonde du \texttt{MockDataProvider} a été menée pour assurer la cohérence de l'application :
\begin{itemize}
    \item \textbf{Résolution des Conflits de Modèles} : Alignement des constructeurs (notamment pour \texttt{Comment}) et des accesseurs (\texttt{getCode} au lieu de \texttt{getInviteCode}) sur les classes POJO réelles.
    \item \textbf{Support Multi-Écrans} : Implémentation des méthodes nécessaires aux activités de notifications, réglages et gestion des membres de groupes.
    \item \textbf{Flexibilité des Données} : Ajout de surcharges de méthodes pour supporter différents types d'appels (par objet ou par identifiant) sans duplication de code.
\end{itemize}

En attendant les endpoints finaux, les données sont fournies par cette architecture locale évolutive.

\subsection{Phase 4 : Affichage Cartographique (\texttt{MapFragment})}
Afin de visualiser géographiquement les photos partagées, un espace cartographique a été intégré via l'onglet "Carte". Ce module repose sur la bibliothèque open-source \textbf{OSMDroid} au lieu du SDK Google Maps, permettant un rendu fidèle à la maquette de style pastel (OpenStreetMap) sans nécessité immédiate de clé API facturable.
Les marqueurs sur la carte sont générés de manière asynchrone en convertissant des layouts XML complexes (contenant les miniatures des photos, un contour blanc et un pointeur de géolocalisation) en objets \texttt{Bitmap} que le moteur de la carte utilise nativement comme icônes d'épingles.

\section{Détails techniques}
\subsection{Stack et Intégrations}
\begin{itemize}
    \item \textbf{Langage} : Java / XML
    \item \textbf{Bases de données} : Cloud Firestore
    \item \textbf{Authentification} : Firebase Authentication
    \item \textbf{Stockage Média} : SDK Cloudinary Android 
    \item \textbf{Cartographie} : OSMDroid (OpenStreetMap)
    \item \textbf{Interface Utilisateur} : ConstraintLayout, Material Design 3.
\end{itemize}

\section{Diagramme de navigation entre les écrans}
(À générer avec TikZ une fois le NavGraph finalisé).

\section{Difficultés rencontrées et solutions apportées}

\subsection{Refonte UI et intégration des VectorDrawables}
Lors de l'implémentation de la charte graphique moderne (icônes "Lucide"), les conversions directes de SVG en \texttt{VectorDrawable} Android ont provoqué des anomalies visuelles complexes sur l'émulateur (icônes partiellement rognées, éléments creux remplis de noir).
\begin{itemize}
    \item \textbf{Problème identifié} : Le parseur graphique d'Android (Skia) gère de manière stricte les coordonnées des \textit{paths} SVG et le remplissage implicite (absence de \texttt{fillColor}).
    \item \textbf{Pistes explorées} : Nous avons intégré les versions \texttt{VectorDrawable} pré-converties issues du projet open-source \textit{0xJihan} (dédié à Android), configuré nos propres teintes via l'attribut \texttt{strokeColor}, et tenté de forcer la recompilation pour invalider le cache d'Android Studio (\texttt{gradle clean}). Les fichiers écrasés par erreur lors de ces tests ont également été restaurés automatiquement ($ic\_launcher$, $chip\_bg$, etc.).
    \item \textbf{État actuel et résolution} : Il s'est avéré que ce problème d'affichage était exclusif au moteur de rendu graphique de l'émulateur Android Studio (artefacts et clipping). Les tests de déploiement réalisés physiquement sur un véritable smartphone Android ont confirmé que les icônes \texttt{VectorDrawable} s'affichent de façon parfaite et sans aucun bug de remplissage, validant ainsi notre intégration native.
\end{itemize}

\subsection{Suppression du thème sombre et verrouillage en mode clair}
Lors des premiers tests sur téléphone physique, l'application héritait automatiquement du mode sombre du système Android, inversant certaines couleurs de la maquette. Pour garantir un rendu fidèle à la charte graphique :
\begin{itemize}
    \item \textbf{Manifest} : Ajout de \texttt{android:forceDarkAllowed="false"} dans la balise \texttt{<application>} du fichier \texttt{AndroidManifest.xml} afin de bloquer la conversion automatique en mode sombre par le système.
    \item \textbf{Thème} : Remplacement du parent \texttt{Theme.Material3.Light.NoActionBar} par \texttt{Theme.MaterialComponents.Light.NoActionBar} pour un comportement plus prévisible sur tous les appareils.
    \item \textbf{Suppression de \texttt{values-night}} : Le dossier \texttt{res/values-night/} a été intégralement supprimé pour éliminer toute surcharge de style nocturne.
    \item \textbf{Attributs dynamiques} : Toutes les références telles que \texttt{?attr/colorSurface} dans les layouts ont été remplacées par des couleurs fixes (\texttt{@android:color/white}) pour éviter toute ambiguïté entre thèmes.
\end{itemize}

\subsection{Gestion de l'encoche (caméra frontale) et des barres système}
Sur les téléphones modernes dotés d'une encoche (``notch'') ou d'un poinçon caméra (``punch-hole''), le contenu de l'application se retrouvait partiellement caché sous la barre d'état système.
\begin{itemize}
    \item \textbf{Cutout Mode} : Ajout de l'attribut \texttt{android:windowLayoutInDisplayCutoutMode="shortEdges"} dans l'activité principale du Manifest pour autoriser le dessin de l'interface jusque sous les bords courts de l'écran.
    \item \textbf{WindowInsetsCompat} : Implémentation d'un listener \texttt{ViewCompat.setOnApplyWindowInsetsListener} dans \texttt{MainActivity.java} qui applique dynamiquement un \texttt{paddingTop} égal à la hauteur de la barre d'état. Le contenu est ainsi décalé juste en-dessous de la caméra frontale sans jamais être masqué.
\end{itemize}

\subsection{Rendu des marqueurs personnalisés sur la carte}
Les marqueurs de la carte (``épingles photo'') s'affichaient comme des carrés bleus unis au lieu du design attendu (vignette photo encadrée de blanc). Ce problème était lié à la conversion hors-écran du composant \texttt{MaterialCardView} en \texttt{Bitmap} : l'élévation et les ombres portées ne sont pas rendues correctement par le moteur \texttt{Canvas} d'Android lorsque la vue n'est pas attachée à une fenêtre.
\begin{itemize}
    \item \textbf{Solution} : Remplacement du layout \texttt{MaterialCardView} par un \texttt{RelativeLayout} utilisant un \texttt{ShapeDrawable} plat (\texttt{bg\_rounded\_white.xml}) et des dimensions de mesure explicites via \texttt{MeasureSpec.EXACTLY}, garantissant un rendu bitmap fidèle même en hors-écran.
\end{itemize}

\subsection{Phase 5 : Interface de Préférences TravelPath}
L'écran ``Créer un parcours'' (\texttt{TravelPathPreferencesFragment}) a été intégré, accessible via l'onglet \textit{Parcours} de la barre de navigation. Il s'agit d'un formulaire défilable (\texttt{ScrollView}) composé de 7 cartes thématiques permettant de configurer un itinéraire :
\begin{itemize}
    \item \textbf{Bannière ``Générer depuis vos photos''} : Un encart à fond dégradé turquoise invitant à créer un parcours automatiquement à partir des photos favorites.
    \item \textbf{Type d'activités} : 4 cases à cocher (Restauration, Loisirs, Découverte, Culture) avec icônes Lucide personnalisées.
    \item \textbf{Lieux favoris} : Chips turquoise pré-remplies (Tour Eiffel, Musée du Louvre, Notre-Dame) avec possibilité d'en ajouter.
    \item \textbf{Budget maximum} : Un \texttt{SeekBar} de 0\euro{} à 500\euro{} avec affichage dynamique de la valeur sélectionnée.
    \item \textbf{Durée} : Sélection exclusive parmi ``½ journée'', ``1 jour'', ``2 jours+''.
    \item \textbf{Profil / Effort} : Sélection exclusive parmi Famille, Seniors, Sportif, Aventurier.
    \item \textbf{Préférences météo} : Chips à bascule (multi-sélection) pour Froid, Chaleur, Humidité, Ensoleillé.
\end{itemize}
Le bouton ``Générer les parcours'' en bas de page servira de point d'entrée vers l'algorithme de génération d'itinéraires.

\subsection{Phase 6 : Vue Recherche Avancée (\texttt{AdvancedSearchFragment})}
L'écran de ``Recherche avancée'' (accessible via l'onglet \textit{Recherche}) a été fidèlement intégré d'après les maquettes. Cette vue complexe est divisée entre des éléments fixes (barre de recherche, sélecteurs d'onglets, et le bouton de validation principal) et un contenu central défilant.
\begin{itemize}
    \item \textbf{En-tête et barre de recherche} : Une barre de recherche au style \textit{pill} (bordures très arrondies, sans contour apparent) flanquée d'un bouton ``Micro'' en carré arrondi.
    \item \textbf{Menu de navigation (Tabs)} : Deux onglets statiques ``Filtres'' (actif, souligné en cyan) et ``Navigation''.
    \item \textbf{Filtres par Type de lieu} : Une grille dynamique affichant les catégories principales (Nature, Musée, Rue, Magasin, Restaurant, Monument). Les pictogrammes de la maquette ont été respectés en utilisant directement les versions vectorielles (SVG converties en \texttt{VectorDrawable}) des icônes \textbf{Lucide}.
    \item \textbf{Filtres complémentaires} : 
    \begin{itemize}
        \item Des blocs sélectionnables en forme de \textit{chips} stylisées pour le "Moment de la journée".
        \item Un menu déroulant masqué pour la "Période".
        \item Des champs textes de saisie avec arrière-plan gris clair pour filtrer par Lieu, Auteur, ou Tags additionnels.
        \item Un bouton spécifique pour rechercher "Autour de ma position" utilisant un \texttt{TextView} stylisé en bouton.
    \end{itemize}
    \item \textbf{Bouton de validation fixe} : Le grand bouton cyan "Rechercher", comportant l'icône loupe, est positionné de façon statique en bas d'écran, restant toujours accessible sans défilement de la page.
\end{itemize}
Le bouton retour de la top-bar a également été interconnecté au \texttt{NavController} sous-jacent via un appel à \texttt{onBackPressed()} de l'activité parente.

\subsection{Phase 7 : Profil Utilisateur (\texttt{ProfileFragment})}
L'écran de profil (Sophie Martin) a été implémenté pour refléter fidèlement la maquette, en incluant les fonctionnalités suivantes :
\begin{itemize}
    \item \textbf{En-tête et statistiques} : Affichage du nom d'utilisateur avec un menu déroulant simulé, et une barre d'actions comprenant l'ajout de contenu (+) et les paramètres (menu hamburger). Les statistiques (publications, abonnés, abonnements) sont présentées en colonnes sous l'avatar.
    \item \textbf{Avatar dynamique} : L'image de profil est affichée de manière circulaire avec un badge flottant "Caméra" (icône Lucide blanche sur fond cyan) permettant de simuler la modification de la photo.
    \item \textbf{Bio et informations} : Affichage du nombre de pays visités avec une icône globe, et d'une courte biographie textuelle.
    \item \textbf{Boutons d'action} : Deux boutons principaux "Modifier le profil" et "Partager le profil" ont été stylisés avec des coins arrondis et un fond gris clair neutre.
    \item \textbf{Onglets et Grille de photos} : Un système d'onglets (Photos / Parcours) permet de basculer l'affichage. L'onglet Photos affiche une grille de 3 colonnes (\texttt{RecyclerView} avec \texttt{GridLayoutManager}) présentant les miniatures des publications de l'utilisateur.
\end{itemize}
Techniquement, les données sont fournies par le \texttt{MockDataProvider} via des méthodes dédiées (\texttt{getCurrentUser} et \texttt{getUserPhotos}). L'interface a été conçue pour forcer visuellement le "mode connecté" afin de présenter l'ensemble des options de gestion de profil, en attendant l'intégration du système d'authentification réel.

\subsection{Phase 11 : Passerelle Inteligente (\texttt{GatewayActivity})}
L'activité \texttt{GatewayActivity} a fait l'objet d'une refonte UI minutieuse pour garantir une fidélité absolue aux maquettes de référence et une expérience utilisateur sans couture.
\begin{itemize}
    \item \textbf{Cohérence Visuelle et Charte Graphique} : Le titre de la Toolbar a été étendu à "Créer un parcours depuis vos photos". L'ensemble de l'iconographie IA a été uniformisée avec l'icône Lucide \texttt{sparkles}. La palette de couleurs a été strictement recalée sur le Cyan (\#0891B2) pour les états actifs, remplaçant l'ancien Teal (\#009688) incorrect.
    \item \textbf{Textes dynamiques et retour à la ligne} : Un algorithme natif \texttt{breakStrategy="simple"} et un width 0dp ont été appliqués aux sous-titres, et le bouton "Générer" autorise désormais son texte sur 2 lignes avec contrainte de `minHeight`, empêchant toute coupure ou écrasement de police.
    \item \textbf{Liste horizontale sans compromis} : Les photos aimées s'affichent correctement en carrousel natif (`LinearLayoutManager.HORIZONTAL`) de 140x180dp sans barre de défilement visuelle parasite (`android:scrollbars="none"`).
    \item \textbf{Alignement parfait des icônes} : La génération programmatique complexe des "Chips" (tags lieux) a été remplacée par un layout dédié \texttt{item\_chip\_place.xml}, offrant un alignement intra-composant rigoureux (icon start padding, text paddings) et un rendu de l'icône \texttt{map-pin} identique à la maquette.
\end{itemize}
Cette phase de finition méticuleuse a permis de valider l'intégration des composants Material Design 3 dans des flux de navigation complexes, tout en assurant une réactivité graphique optimale.

\subsection{Phase 8 : Panneau de Notifications (\texttt{NotificationsFragment})}
L'écran de notifications a été implémenté en tant que fragment complet, accessible via l'icône cloche de l'accueil. Cette vue permet de centraliser les interactions sociales et les alertes système :
\begin{itemize}
    \item \textbf{En-tête et actions} : Une barre supérieure fixe affichant le nombre de notifications non lues, un bouton d'accès aux paramètres et une option "Tout marquer" pour valider l'ensemble des notifications d'un seul clic.
    \item \textbf{Système d'onglets} : Deux onglets dynamiques ("Toutes" et "Non lues") permettent de filtrer les notifications par état. Chaque onglet affiche un badge numérique synchronisé avec les données.
    \item \textbf{Liste de notifications riche} : Le \texttt{RecyclerView} utilise l'adaptateur \texttt{NotificationAdapter} pour gérer différents types des contenus :
    \begin{itemize}
        \item Interactions sociales (Likes, Commentaires, Abonnements) avec l'avatar de l'auteur et, si applicable, une miniature de la photo concernée.
        \item Alertes système (Itinéraires prêts, Suggestions de lieux) avec des icônes thématiques dédiées sur fond coloré.
    \end{itemize}
    \item \textbf{États visuels} : Les notifications non lues sont distinguées par un point bleu et un arrière-plan très légèrement contrasté. Cliquer sur un onglet ou sur "Tout marquer" met à jour dynamiquement l'interface.
\end{itemize}
Le \texttt{NotificationsFragment} est alimenté par la méthode \texttt{getNotifications} de \texttt{MockDataProvider}, simulant un historique varié de 8 notifications. L'intégration avec Firebase Cloud Messaging (FCM) est prévue pour le passage en production.

\subsection{Phase 9 : Paramètres des notifications (\texttt{NotificationSettingsActivity})}
Cet écran permet à l'utilisateur de fine-tuner ses alertes par catégorie (personnes, groupes, lieux, tags). Conformément à la maquette :
\begin{itemize}
    \item \textbf{Contrôle Global} : Un switch principal permet d'activer ou désactiver l'ensemble des notifications push via une carte mise en évidence.
    \item \textbf{Gestion granulaire} : Quatre sections distinctes permettent de gérer les abonnements spécifiques. Chaque élément est présenté dans une carte blanche avec son propre switch d'activation et une option de suppression brute.
    \item \textbf{Architecture dynamique} : L'activité génère dynamiquement les vues à partir des données fournies par \texttt{MockDataProvider.getNotificationSettings()}. Les interactions (basculement, suppression) sont gérées localement en mémoire pour simuler le comportement final.
    \item \textbf{Design System} : Utilisation du code couleur Teal pour les Switchs et les boutons d'ajout. La section "Tags" utilise un rendu spécifique sous forme de Chips pour une meilleure distinction visuelle.
\end{itemize}
Les classes \texttt{NotificationSettingsActivity} (Logique UI) et \texttt{NotificationSettingItem} (Modèle de données) constituent le cœur de ce module. Les données sont actuellement isolées de Firebase pour cette phase de prototypage UI/UX.

\subsection{Phase 10 : Gestion des Groupes (\texttt{GroupsActivity})}
Cette phase introduit la gestion des communautés au sein de l'application. L'architecture repose sur un système d'onglets pour séparer l'expérience utilisateur :
\begin{itemize}
    \item \textbf{Navigation par onglets} : Utilisation de \texttt{TabLayout} synchronisé avec un \texttt{ViewPager2}. Deux fragments \texttt{GroupListFragment} sont instanciés pour gérer respectivement "Mes groupes" et "Découvrir".
    \item \textbf{Cartes de Groupes} : Un \texttt{GroupAdapter} gère le rendu complexe des cartes. Chaque carte affiche un badge dynamique (Public/Privé en Teal/Gris), le compteur de membres avec l'icône \texttt{Users}, le compteur de photos avec l'icône \texttt{Camera}, et une icône \texttt{Crown} (Couronne) exclusive aux administrateurs.
    \item \textbf{Modèle et Données} : La classe \texttt{Group} encapsule les propriétés sociales. Les données sont fournies par \texttt{MockDataProvider.getMyGroups()} et \texttt{getDiscoverGroups()}, simulant un environnement peuplé.
    \item \textbf{Transitions} : Le clic sur un groupe navigue vers \texttt{GroupDetailActivity} (stub) ou ouvre le dialogue de détail riche selon le contexte utilisateur.
\end{itemize}
Techniquement, la persistance des données pour ce module a été refactorisée. \texttt{MockDataProvider} utilise désormais des listes statiques persistantes en mémoire vive. Ce choix permet d'ajouter dynamiquement des groupes (via création ou adhésion) et de voir ces changements se répercuter immédiatement sur l'ensemble des fragments, simulant le comportement d'une base de données locale ou distante.
L'accès se fait directement depuis la barre d'outils de l'écran d'accueil via le bouton \texttt{ivGroup}.

\subsubsection{Composant : \texttt{JoinGroupDialogFragment}}
Pour permettre l'adhésion à des groupes restreints, un fragment de dialogue a été implémenté :
\begin{itemize}
    \item \textbf{Interface} : Un dialogue aux coins très arrondis (20dp) conforme à la charte Material Design 3, incluant un champ de saisie stylisé en Cyan.
    \item \textbf{Validation} : Le bouton "Rejoindre" est activé dynamiquement par un \texttt{TextWatcher} dès qu'un caractère est saisi. Le texte est automatiquement transformé en majuscules pour correspondre au format des codes d'invitation.
    \item \textbf{Affinage UI} : Les boutons "Annuler" et "Rejoindre" ont été calibrés avec une taille de texte de 16sp et la désactivation de la mise en majuscules automatique (\texttt{textAllCaps=false}) pour garantir une lisibilité optimale et une fidélité accrue au design original.
    \item \textbf{Logique Métier} : La vérification s'appuie sur \texttt{MockDataProvider.findGroupByCode()}. Les codes \texttt{PAR2026} (Voyage Paris) et \texttt{FAM-MTN} (Famille Martin) sont reconnus. En cas d'échec, un message d'erreur rouge s'affiche dynamiquement sous le champ.
\end{itemize}
Ce composant sera ultérieurement relié à une fonction Firebase pour valider les tokens d'invitation en temps réel.

\subsubsection{Composant : \texttt{CreateGroupDialogFragment}}
Complémentaire au module de ralliement, ce dialogue permet l'initialisation de nouvelles communautés :
\begin{itemize}
    \item \textbf{Gestion de la Visibilité} : Mise en œuvre d'un sélecteur de visibilité agissant comme des boutons radio riches. L'option sélectionnée (Public ou Privé) bénéficie d'une bordure Cyan et d'un fond teinté (\texttt{bg\_radio\_option\_active}), tandis que l'option inactive repasse en mode neutre.
    \item \textbf{Génération de Code} : Pour les groupes marqués comme "Privé", un algorithme de génération aléatoire produit un code alphanumérique de 6 caractères (ex: \texttt{X7B92A}). Ce code est immédiatement communiqué à l'utilisateur via un Toast persistant.
    \item \textbf{Flux de Données} : Le nouveau groupe est inséré dynamiquement en tête de la liste \texttt{myGroups} via \texttt{MockDataProvider.addGroup()}. Un mécanisme d'écouteur (\texttt{OnGroupCreatedListener}) permet à l'activité parente de rafraîchir instantanément le \texttt{ViewPager2}.
\end{itemize}
La persistance réelle des données sera assurée par l'intégration de Firebase Firestore dans une phase ultérieure.

\subsubsection{Correction : Navigation Inter-Activités}
Un bug identifié dans la barre de navigation de \texttt{GroupsActivity} a été corrigé. Bien que visuellement présente, la barre ne permettait pas de naviguer vers les autres sections de l'application. La solution implémentée utilise des \texttt{Intent} avec des drapeaux de type \texttt{FLAG\_ACTIVITY\_CLEAR\_TOP} pour réutiliser l'instance de \texttt{MainActivity} et y injecter l'ID du fragment cible via un extra \texttt{target\_fragment\_id}. \texttt{MainActivity} a été modifiée pour intercepter ces appels dans \texttt{onNewIntent} et déclencher la navigation via son \texttt{NavController}. Un crash lié à une erreur d'ID lors de l'initialisation de cette barre (\texttt{NullPointerException}) a également été résolu par une synchronisation rigoureuse des identifiants XML et Java.

\subsubsection{Composant : \texttt{GroupDetailDialogFragment}}
Une refonte majeure de l'écran de détail a été effectuée, passant d'une activité standard à un \texttt{DialogFragment} riche et contextuel. Ce composant introduit une gestion fine des droits via l'énumération \texttt{UserRole} :
\begin{itemize}
    \item \textbf{ADMIN} : Accès total. Affiche le bandeau narratif, les statistiques complètes, le bloc de code d'invitation avec fonction de copie, ainsi que le bouton "Gérer".
    \item \textbf{MEMBER\_WITH\_CODE} : Rôle de membre ayant le droit de partager le groupe. Affiche le code d'invitation mais masque le bouton de gestion administrative.
    \item \textbf{MEMBER} : Accès restreint à la consultation. Masque totalement le bloc de code d'invitation pour garantir la confidentialité des accès privés.
\end{itemize}
Techniquement, l'ouverture est déclenchée par le \texttt{GroupAdapter} via l'injection d'un \texttt{FragmentManager}. Le design inclut un système d'avatars dynamiques générés à partir des initiales des membres, avec un badge de débordement (\texttt{+N}) pour les groupes importants.

\subsubsection{Onglet : Découverte de Groupes}
L'onglet "Découvrir" a été implémenté via le fragment \texttt{DiscoverGroupsFragment}. Il propose une expérience de recherche immersive avec une barre de filtrage en temps réel.
\begin{itemize}
    \item \textbf{Filtrage Dynamique} : L'adaptateur \texttt{DiscoverGroupAdapter} implémente un mécanisme de filtrage basé sur le nom et la description des groupes.
    \item \textbf{Design Distingué} : Contrairement aux groupes personnels, les cartes de découverte utilisent \texttt{item\_group\_discover.xml} avec une bannière pleine largeur et un overlay en gradient pour le nom du groupe.
    \item \textbf{Logique d'Adhésion} : Le bouton d'état gère visuellement le passage de "Rejoindre" (Teal) à "Rejoint" (Gris/Désactivé). L'adhésion déclenche une mise à jour dans \texttt{MockDataProvider}, ajoutant instantanément le groupe à l'onglet "Mes groupes" avec le rôle \texttt{MEMBER}.
\end{itemize}
Cette architecture permet une séparation nette entre la consultation des groupes rejoints et l'exploration de nouvelles communautés.

\subsection{Phase 12 : Écran de Paramètres globaux (\texttt{SettingsActivity})}
La gestion du compte et des préférences de l'utilisateur a été matérialisée par la création from scratch de l'écran \texttt{SettingsActivity}, accessible depuis l'icône \texttt{menu} (Lucide) du \texttt{ProfileFragment}.
Cette activité structure proprement les différentes options selon la maquette fournie, en listant les sections via un \texttt{NestedScrollView} pour un défilement complet.

\begin{itemize}
    \item \textbf{Carte Utilisateur (\texttt{MaterialCardView})} : Affiche les informations de l'utilisateur connecté via chargement dynamique depuis le \texttt{MockDataProvider}.
    \item \textbf{Section COMPTE} :
    \begin{itemize}
        \item \textit{Modifier le profil} : Débouche sur \texttt{EditProfileActivity} (stub actuel).
        \item \textit{Notifications} : Redirige vers la \texttt{NotificationSettingsActivity}.
    \end{itemize}
    \item \textbf{Section PRÉFÉRENCES} :
    \begin{itemize}
        \item \textit{Langue} : Ouvre un `AlertDialog` gérant le basculement fonctionnel simulé de la sélection (Français, English, Español).
        \item \textit{Confidentialité} : Débouche sur \texttt{PrivacyActivity} (stub actuel).
    \end{itemize}
    \item \textbf{Section SUPPORT} :
    \begin{itemize}
        \item \textit{Aide \& Support} : Débouche sur \texttt{SupportActivity} (stub actuel).
    \end{itemize}
    \item \textbf{Déconnexion} : Un bloc cliquable rouge clair distinct du reste de la liste lance une demande de confirmation via dialogue. Actuellement, l'action `finishAffinity()` purge l'historique et redirige vers \texttt{MainActivity}.
\end{itemize}

\textbf{Information liée aux données et futures évolutions :} 
L'ensemble des options et profils sont actuellement mockés. Cet écran est structurellement prêt à accueillir \textbf{Firebase Auth} (pour la structure dynamique du profile et l'action de déconnexion) ainsi que \textbf{Firestore} pour l'enregistrement persistant des préférences utilisateur (langue, confidentialité). Lors du raccordement au Backend, le bouton "Se déconnecter" pointera logiquement vers une \texttt{LoginActivity}.

\subsubsection{Correctifs UI : Insets et Icônes}
Une passe de stabilisation a corrigé des anomalies visuelles urgentes :
\begin{itemize}
    \item \textbf{Harmonisation des teintes} : Les icônes "menu" et "ajouter" du \texttt{ProfileFragment} se sont vu attribuer une dynamique via \texttt{app:tint="@color/on\_background"} pour matcher la couleur des autres ressources (claire ou sombre selon le contexte), remplaçant l'affichage blanc cassé.
    \item \textbf{Gestion des Marges (Edge-To-Edge)} : Suppression d'effets de double-marges (grand écart en haut de l'écran) identifiés dans le \texttt{ProfileFragment}. La redondance entre l'attribut \texttt{fitsSystemWindows="true"} du layout XML et l'application manuelle des Insets (\texttt{ViewCompat.setOnApplyWindowInsetsListener}) a été révoquée. L'approche s'en remet désormais au \texttt{MainActivity} pour le fragment. En revanche, \texttt{SettingsActivity} étant une activité autonome (hors du NavController), son propre \texttt{ViewCompat.setOnApplyWindowInsetsListener} lui a été rétabli de manière exclusive afin d'éviter qu'elle ne passe sous la caméra frontale, garantissant un flux Inset naturel et homogène.
\end{itemize}

\subsection{Phase 13 : Édition du Profil (\texttt{EditProfileActivity})}
Afin de compléter la gestion de compte, l'écran \texttt{EditProfileActivity} a été implémenté d'après les maquettes, avec une structure basée sur un \texttt{CoordinatorLayout} et un formulaire riche au sein d'un \texttt{NestedScrollView}. Elle s'ouvre logiquement depuis l'item "Modifier le profil" de l'écran des paramètres.

\subsubsection{Interface et Formulaire}
\begin{itemize}
    \item \textbf{Zone Avatar} : Conteneur \texttt{ShapeableImageView} circulaire doté d'un "badge" caméra statique (overlap). Clyn d'œil interactif : un clic déclenche l'\texttt{Intent.ACTION\_PICK} pour sélectionner une image locale et met à jour dynamiquement l'ImageView.
    \item \textbf{Champs configurés} : Le formulaire intègre des \texttt{EditText} stylisés globalement via un sélecteur XML (bordure Cyan `#0891b2` au focus). Les champs intègrent Nom, Localisation (avec icône intra-champ), Email et Site Web.
    \item \textbf{Compteur biographique} : Le champ de biographie, restreint à \texttt{150 caractères} (\texttt{maxLines=5}), est observé par un \texttt{TextWatcher}. Celui-ci met visuellement à jour le label de progression ("X / 150 caractères").
    \item \textbf{Section Lecture Seule} : Recapitulatif inerte incluant les statistiques basiques (ex: l'ancienneté hardcodée à "Janvier 2026").
\end{itemize}

\subsubsection{Gouvernance et Zone de Danger}
\begin{itemize}
    \item \textbf{Persistance Mocks} : Le \texttt{MockDataProvider} a été refactoré pour cacher l'objet \texttt{currentUser} sous forme de \textit{Singleton}. Cela permet au bouton "Enregistrer" de la \textit{Toolbar} d'éditer réellement ce \texttt{User} mocké et de vérifier la non-vacuité du champ Nom avant fermeture de la page, répercutant l'information aux écrans parents.
    \item \textbf{Modales Destructives} : Les boutons de désactivation et de suppression du compte activent des \texttt{AlertDialog} d'alerte, peints avec l'accentuation `#EF4444` (Rouge). Une confirmation purge virtuellement l'application (\texttt{finishAffinity()}) dans sa version actuelle.
\end{itemize}

\textbf{Information liée aux données et futures évolutions :} 
Toutes les saisies et interactions media (\texttt{URI}) s'arrêtent actuellement au cache de l'appareil. Le système est structuré pour une connexion ultérieure à \textbf{Firebase Storage} pour l'upload de l'avatar et à \textbf{Firestore} pour l'enregistrement robuste des champs de profil et du rôle d'authentification.

\subsection{Phase 14 : Partage de Profil et Parcours Générés}

\subsubsection{Composant : \texttt{ShareProfileDialogFragment}}
Le bouton "Partager le profil" ouvre désormais une boîte de dialogue contextuelle (non pas un BottomSheet, mais un dialogue centré à bords arrondis), facilitant l'export de la fiche utilisateur vers l'écosystème du téléphone.
\begin{itemize}
    \item \textbf{Design} : Titre "Partager" flanqué de l'icône \texttt{share-2} et d'un bouton de fermeture. La liste d'options est stylisée de manière aérée avec des icônes dans des cercles gris clair pour un rendu optimal.
    \item \textbf{Interactions Systèmes} :
    \begin{itemize}
        \item \textit{Partage Système} : Exploitation de \texttt{Intent.ACTION\_SEND} avec un type MIME \texttt{text/plain} pour ouvrir le tiroir de partage natif d'Android.
        \item \textit{Presse-papiers} : Utilisation du service \texttt{ClipboardManager} pour copier discrètement l'URL d'invitation avec confirmation via \texttt{Toast}.
        \item \textit{Email \& SMS} : Appels directs aux Intents \texttt{ACTION\_SENDTO} avec les préfixes \texttt{mailto:} et \texttt{smsto:} pour injecter un texte de découverte préparé à l'avance.
    \end{itemize}
    \item \textbf{Mocking de l'URL} : En attendant l'intégration des Dynamic Links de Firebase, les URL partagées reprennent le nom du \texttt{User} mocké (ex: \texttt{https://traveling.app/u/sophie.martin}).
\end{itemize}

\subsubsection{Écran : \texttt{GeneratedItinerariesActivity}}
Le bouton de génération au sein de \texttt{TravelPathPreferencesFragment} aboutit désormais sur une nouvelle activité présentant les parcours concoctés par le système, \texttt{GeneratedItinerariesActivity}.
\begin{itemize}
    \item \textbf{Modèle de Données} : Création de la classe \texttt{GeneratedItinerary} stockant les métriques principales (Titre, Description, Budget, Durée, Niveau d'effort, Nombre d'étapes) ainsi que l'état d'appréciation ("Like"). Trois profils factices (Économique, Équilibré, Confort) sont renvoyés par \texttt{MockDataProvider}.
    \item \textbf{Bannière Contextuelle} : Un encart informatif (bleu \texttt{\#E0F7FA}) au-dessus de la liste rassure l'utilisateur en synthétisant avec mise en gras \texttt{Html.fromHtml()} le nombre d'itinéraires trouvés.
    \item \textbf{Carte Itinéraire} : Implémentée avec le \texttt{MaterialCardView}, chaque recommandation présente une matrice $2 \times 2$ très lisible pour les caractéristiques du parcours.
    \item \textbf{Badges Dynamiques} : Le niveau d'effort est soutenu visuellement par un \texttt{GradientDrawable} dont la couleur change de manière programmatique selon la sévérité (Vert clair pour "Facile", Jaune pastel pour "Modéré", et Rouge rosé pour "Difficile").
    \item \textbf{Bouton Like Actif} : L'action sur l'icône de cœur bascule non seulement le booléen du modèle, mais redessine l'interface avec la version "pleine" de l'icône et teinte le tracé en \texttt{\#EF4444} (Rouge).
\end{itemize}
Dans les prochaines itérations, ces données seront générées par le backend IA basé sur les choix formulés dans le formulaire précédent. Un écran final \texttt{ItineraryDetailActivity} a été stubs pour recueillir le détail de la navigation point-par-point plus tard.

\section{Phase 15 : Détails de Parcours, Export PDF et Profil Utilisateur}
\subsection{Description}
Cette phase a permis d'implémenter l'écran complet de détail d'un parcours généré par l'IA, la modale d'export PDF associée, ainsi que le peuplement de l'onglet "Parcours" sur le profil de l'utilisateur.

\subsection{Détails techniques}
\begin{itemize}
    \item \textbf{ItineraryDetailActivity} : Une activité complexe utilisant un \texttt{NestedScrollView}. Elle est structurée en 5 blocs distincts :
    \begin{itemize}
        \item Bloc 1 : Carte statique et bouton flottant redirigeant vers une application de navigation GPS externe via Intent (\texttt{geo:}).
        \item Bloc 2 : Météo du jour présentée dans une grille 2x2 avec icônes personnalisées.
        \item Bloc 3 : Statistiques résumées (budget, durée, nombre d'étapes).
        \item Bloc 4 : Option de sauvegarde hors-ligne simulant un téléchargement local.
        \item Bloc 5 : Une "Timeline" générée dynamiquement affichant la chronologie des étapes via \texttt{item\_itinerary\_step.xml}.
    \end{itemize}
    \item \textbf{Modèles et Données} : Création des classes \texttt{ItineraryStep.java} et \texttt{ProfileItinerary.java}. Le \texttt{MockDataProvider} a été enrichi pour fournir les données associées.
    \item \textbf{Export PDF} : \texttt{PdfReadyDialogFragment} est un \texttt{DialogFragment} centré simulant la confirmation de génération PDF. Il inclut une fonctionnalité de partage via \texttt{Intent.ACTION\_SEND} avec le type MIME \texttt{application/pdf}.
    \item \textbf{Onglet Parcours du Profil} : Le \texttt{ProfileFragment} intègre désormais un second \texttt{RecyclerView} (\texttt{rvProfileTrips}) alimenté par \texttt{ProfileItineraryAdapter}. Le basculement entre l'onglet "Photos" et "Parcours" gère dynamiquement la visibilité des deux listes, en plus de modifier l'état actif des boutons.
\end{itemize}

\subsection{Difficultés et Solutions}
\begin{itemize}
    \item \textbf{Difficulté} : La quantité d'icônes spécifiques (Météo, Actions) nécessaires pour l'écran de détail du parcours.
    \item \textbf{Solution} : Création de 10 nouveaux \texttt{VectorDrawable} (Lucide Icons) directement en XML (\texttt{ic\_sun}, \texttt{ic\_cloud}, \texttt{ic\_droplets}, \texttt{ic\_wind}, \texttt{ic\_thermometer}, etc.) pour maintenir le poids de l'application bas et garantir une mise à l'échelle parfaite, sans ajouter de dépendances externes.
\end{itemize}

\subsection{Améliorations UI/UX (Correctifs)}
Suite à une revue détaillée, plusieurs améliorations esthétiques et fonctionnelles ont été apportées :
\begin{itemize}
    \item \textbf{Lisibilité des Blocs} : Les encarts "Météo" et "Statistiques" bénéficient de nouveaux fonds dédiés (cyan très pâle \texttt{\#ECFCFF} et \texttt{\#F2F9FB}) accompagnés de bordures fines afin de structurer l'information visuellement.
    \item \textbf{Galerie Média par Étape} : Remplacement du layout simple par une vue déroulante (Expand/Collapse). Un composant interactif affiche désormais un \texttt{RecyclerView} horizontal de miniatures (via le nouveau modèle \texttt{StepPhoto.java}) lorsqu'on clique sur le bandeau média d'une étape. Au clic sur une miniature, une modale plein écran (\texttt{MediaViewerDialogFragment}) s'ouvre pour afficher l'image ou lire la vidéo, offrant une expérience immersive.
    \item \textbf{Ajustement Modale PDF} : La largeur du \texttt{PdfReadyDialogFragment} a été fixée dynamiquement via \texttt{onStart()} à 90\% de la largeur de l'écran (\texttt{widthPixels * 0.90}) pour éviter l'écrasement du contenu.
    \item \textbf{Mutualisation du Partage} : Le \texttt{ShareProfileDialogFragment} a été rendu réutilisable grâce à une méthode \texttt{newInstanceForItinerary} injectant dynamiquement via un \texttt{Bundle} les textes et URLs adaptés au contexte d'un parcours (au lieu d'un profil).
    \item \textbf{Correction des Badges} : Application de \texttt{singleLine="true"} pour empêcher le retour à la ligne du créneau "Après-midi", et ajustement de la colorimétrie des fonds (notamment le violet pâle \texttt{\#EDE7F6} pour le badge "Soir") pour garantir un contraste optimal.
\end{itemize}

\subsection{Phase 16 : Intégration de la Navigation Globale}
L'un des défis majeurs de l'application résidait dans la coexistence d'une architecture à Fragment unique (\texttt{MainActivity}) et d'activités autonomes pour les flux complexes (Génération d'itinéraires, Détails). Cette phase a permis d'unifier l'expérience utilisateur en intégrant la barre de navigation inférieure (\texttt{BottomNavigationView}) dans toutes les pages clés.

\subsubsection{Détails techniques}
\begin{itemize}
    \item \textbf{Mutualisation de la Navigation} : La barre de navigation a été ajoutée aux layouts \texttt{activity\_generated\_itineraries.xml}, \texttt{activity\_itinerary\_detail.xml} et \texttt{activity\_gateway.xml}. Pour maintenir une structure propre, ces layouts ont été convertis en \texttt{ConstraintLayout} permettant de fixer le contenu entre la barre d'outils et la barre de navigation.
    \item \textbf{Logique de Routage Inter-Activités} : Comme ces activités ne sont pas gérées par le \texttt{NavController} de l'activité principale, une logique de redirection manuelle a été implémentée dans chaque classe Java (\texttt{setupBottomNavigation}). 
    \item \textbf{Drapeaux de Navigation} : L'utilisation de \texttt{Intent.FLAG\_ACTIVITY\_CLEAR\_TOP} et \texttt{FLAG\_ACTIVITY\_SINGLE\_TOP} garantit que le retour vers l'accueil ou le profil ne crée pas de multiples instances de \texttt{MainActivity}, préservant ainsi la pile d'activités (backstack).
    \item \textbf{Synchronisation des États} : L'ID du fragment cible est passé en extra (\texttt{target\_fragment\_id}), permettant à \texttt{MainActivity} d'ouvrir directement le bon onglet (Accueil, Carte, Parcours, Recherche ou Profil) dès le démarrage.
\end{itemize}

\subsubsection{Difficultés et Solutions}
\begin{itemize}
    \item \textbf{Problème} : Le contenu défilant (\texttt{NestedScrollView}) passait derrière la barre de navigation ou était rogné en bas de page.
    \item \textbf{Solution} : Utilisation des contraintes \texttt{app:layout\_constraintBottom\_toTopOf="@id/bottom\_navigation"} et d'une hauteur \texttt{0dp} pour le conteneur de contenu, forçant Android à calculer dynamiquement l'espace disponible entre les composants fixes.
\end{itemize}

\subsection{Améliorations UI/UX (Affinage Final)}
Suite aux retours utilisateurs, plusieurs ajustements de précision ont été effectués :
\begin{itemize}
    \item \textbf{Nettoyage des Cartes} : Suppression des icônes de favoris (cœurs) redondantes en haut à droite des cartes d'itinéraires générés, pour ne conserver que l'interaction principale en bas à gauche, évitant ainsi la surcharge visuelle.
    \item \textbf{Cohérence des Couleurs} : Recalage des fonds de badges (Cyan \texttt{\#06B6D4} pour la marche, Gris clair \texttt{\#F3F4F6} pour la météo nuageuse) pour une harmonie parfaite avec la charte Tailwind demandée.
    \item \textbf{Correction des Icônes} : L'icône de l'euro (\texttt{ic\_euro}) a été redessinée manuellement en XML pour corriger un défaut d'alignement des barres horizontales.
\end{itemize}

\subsection{Phase 17 : Barre d'actions Contextuelle (Parcours Générés)}
Afin d'améliorer l'interactivité de l'écran des parcours générés, une nouvelle barre d'actions a été intégrée en fin de liste.

\begin{itemize}
    \item \textbf{Interface Sobres} : Contrairement aux boutons Material classiques, ces boutons utilisent un style ``Outline'' avec un fond blanc cassé (\texttt{\#F8FAFB}) et une bordure fine grise. Ce design a été implémenté via un \texttt{drawable} personnalisé (\texttt{bg\_action\_btn\_outline.xml}).
    \item \textbf{Structure et Répartition} : Les deux boutons (``Modifier les préférences'' et ``Regénérer'') occupent chacun 50\% de la largeur de l'écran via un système de \texttt{layout\_weight}.
    \item \textbf{Positionnement Statique} : Conformément aux exigences d'ergonomie, la barre est positionnée à la suite du contenu défilant (\texttt{NestedScrollView}). Elle ne "flotte" pas sur l'interface mais apparaît naturellement à la fin de la liste des itinéraires.
    \item \textbf{Correction des Badges (FlowLayout)} : Pour résoudre définitivement les problèmes de texte tronqué sur les petits écrans, les lignes d'informations des étapes utilisent désormais des \texttt{ChipGroup}. Ce composant permet un retour à la ligne automatique (wrapping) des badges (Ouvert, Matin, etc.) si l'espace horizontal est insuffisant, garantissant l'intégrité visuelle du texte en toute circonstance.
    \item \textbf{Synchronisation de la Navigation} : Correction d'un bug où la barre de navigation basse perdait sa sélection après plusieurs changements de pages. L'implémentation utilise désormais \texttt{setSelectedItemId} dans la \texttt{MainActivity} pour forcer la synchronisation de l'état visuel avec le contrôleur de navigation, et le \texttt{launchMode} a été passé en \texttt{singleTop} pour une gestion plus saine de la pile d'activités.
    \item \textbf{Optimisation de l'Espace} : Réduction des marges de la timeline et passage du bouton de déploiement des photos en pleine largeur pour maximiser l'espace utile sur les écrans étroits.
    \item \textbf{Interactions} :
    \begin{itemize}
        \item Le bouton de modification déclenche un \texttt{finish()}, ramenant naturellement l'utilisateur au formulaire de préférences.
        \item Le bouton de régénération affiche un \texttt{Snackbar} de confirmation, simulant le lancement d'un nouveau calcul par l'IA.
    \end{itemize}
\end{itemize}

\end{document}
